% ============================================================
% BASELINE — Who&When adaptation for the TRAIL benchmark
%
% Two blocks:
%   (1) main-text snippet: 1-2 sentences for the body of the paper
%   (2) appendix block:    full methodological detail
%
% Source paper:
%   Yin et al., "Which Agent Causes Task Failures and When?
%   On Automated Failure Attribution of LLM Multi-Agent Systems"
%   (ICML 2025 Spotlight; arXiv:2505.00212)
%
% Adapted code lives in:
%   baselines/who_and_when/run_who_and_when_vllm.py
% ============================================================


% ---------- (1) MAIN-TEXT SNIPPET (1-2 sentences) -----------------

\paragraph{Who\&When baselines.}
We additionally compare against the all-at-once (\textsc{W1}) and
step-by-step (\textsc{W2}) localization strategies of
\citet{yin2025whoandwhen}, adapted from single-agent attribution to
TRAIL's multi-label, span-level setting. We omit the binary-search
variant (\textsc{W3}): its $\mathcal{O}(\log N)$ efficiency claim
relies on a single-error assumption that is broken on TRAIL, where
the multi-label re-design forces $\mathcal{O}(K\!\cdot\!N)$ leaf
calls in the worst case and offers no compute saving over a
\textsc{W2} sweep.
\textsc{W2}'s $\mathcal{O}(N)$ call schedule also carries a
substantial runtime cost relative to all $\mathcal{O}(1)$ methods:
end-to-end, \textsc{W2} requires 12–23 GPU-hours per model to
process the full dedup evaluation set (234 TRAIL-GAIA +
62 TRAIL-SWE-Bench traces), versus 80–96 minutes for our
\textsc{+GI+SI} method — an 8–17$\times$ wall-time gap at equal or
lower GPU allocation (see Appendix~\ref{app:w2_efficiency}).


% ---------- (2) APPENDIX BLOCK (full detail) ----------------------

\section{Who\&When Adaptation Details}
\label{app:whoandwhen_trail}

\citet{yin2025whoandwhen} introduce three localization strategies
for the failure-attribution problem on multi-agent systems:
\textsc{W1} (all-at-once), \textsc{W2} (step-by-step), and
\textsc{W3} (binary search). All three answer a single
question --- \emph{which agent caused the failure, and at which
step?} --- and assume exactly one responsible agent and one
decisive step per trace, returning a single \texttt{(agent, step)}
tuple. We adapt these strategies to TRAIL by remapping
``\emph{who}'' (which agent) to ``\emph{what}'' (which leaf
category from TRAIL's 19-class taxonomy) and ``\emph{when}'' (which
step) to ``\emph{where}'' (which \texttt{span\_id} in the trace),
while removing the single-error assumption so that the predictor
can return $0$, $1$, or many \texttt{(category, span\_id)} pairs
per trace. The adapted runner is
\texttt{baselines/who\_and\_when/run\_who\_and\_when\_vllm.py};
all variants emit the standard TRAIL output schema so the
existing scorer (\texttt{eval/calculate\_scores.py}) is reused
without modification.

\paragraph{W1 --- All-at-Once.}
The model receives the full trace in a single prompt that mirrors the
original \textsc{W1} prompt structure: ``You are an AI assistant tasked
with analyzing an AI agent execution trace when solving a real-world
problem. The problem is: \{task\_description\}. \ldots Identify which
\{categories\} are present, at which \{span\_id\}, and explain the
reason.'' Compared to \citet{yin2025whoandwhen} we (i) replace the
``which agent caused the failure, at which step'' axis with TRAIL's
19-leaf taxonomy on the category axis and \texttt{span\_id} on the
location axis, (ii) remove the \texttt{errors[0]} truncation so the
output is multi-label, and (iii) drop \texttt{ground\_truth} from the
prompt because TRAIL is reference-free. The task description is kept
verbatim from the original framing. One LLM call per trace.

\paragraph{W2 --- Step-by-Step.}
The original \textsc{W2} walks the trace one step at a time, asks a
yes/no question at each step, and \emph{stops at the first ``yes''}.
We remove the early-exit and replace the per-step prompt with a
multi-label JSON request that may return zero or more categories
per step; otherwise the prompt body, the task description, and the
original's ``\emph{Note: avoid being overly critical \ldots focus on
errors that clearly derail the process}'' calibration sentence are
retained verbatim from \citet{yin2025whoandwhen}. Per-step
predictions are aggregated by deduplicating
\texttt{(category, span\_id)} pairs across the full sweep, with
\texttt{span\_id} validation against the trace's span table. The
sweep covers every step-level span; the only stopping condition is
context overflow --- the original Who\&When relies on a reactive
API-failure bail-out (Azure OpenAI rejects the call), which is fine
for short multi-agent chat histories but unreliable for longer TRAIL
traces, so we add a proactive token-count check that records
\texttt{context\_overflow\_at\_step\_$i$} in the per-trace metadata
and ends the sweep at that point. This is a TRAIL-side safety belt
with the same end-state behaviour as the original, not a methodological
change. $N$ LLM calls per trace, where $N$ is the number of step-level
spans actually scanned (measured median: $N{=}10$ on TRAIL-GAIA,
$N{=}30$ on TRAIL-SWE-Bench, across completed \textsc{W2} traces).

\paragraph{Computational cost of \textsc{W2}.}
\label{app:w2_efficiency}
The $\mathcal{O}(N)$ call schedule of \textsc{W2} produces a
substantial end-to-end runtime gap relative to the $\mathcal{O}(1)$
methods.
Table~\ref{tab:w2_efficiency} reports wall time and GPU-hours for
\textsc{W2} and \textsc{+GI+SI} on the three models for which both
runs are available (Gemma-3-27B-IT, GPT-oss-20B, GPT-oss-120B),
processing all 296 dedup traces (234 TRAIL-GAIA + 62
TRAIL-SWE-Bench).
\textsc{+GI+SI} runtimes are measured from completed Slurm jobs;
\textsc{W2} totals are extrapolated from throughput measured over
multiple 4-hour job windows (the method never finished a single window
for any model, requiring between two and eight sequential reruns per
model-dataset pair).

\begin{table}[h]
\centering
\small
\caption{Wall time and GPU-hours for \textsc{W2} (estimated) vs.\
\textsc{+GI+SI} (measured) on the full 296-trace dedup split.
\textsc{W2} GPU allocations follow the baseline scripts (2 GPUs for
Gemma, 4 GPUs for GPT-oss-20B and GPT-oss-120B);
\textsc{+GI+SI} uses 2 GPUs for Gemma and GPT-oss-20B, 4 for
GPT-oss-120B.}
\label{tab:w2_efficiency}
\begin{tabular}{lrrrr}
\toprule
 & \multicolumn{2}{c}{\textsc{W2} (est.)} & \multicolumn{2}{c}{\textsc{+GI+SI} (measured)} \\
\cmidrule(lr){2-3}\cmidrule(lr){4-5}
Model & Wall time & GPU-hr & Wall time & GPU-hr \\
\midrule
Gemma-3-27B-IT    & ${\sim}12$\,h & 23.9 & 91\,min & 3.0 \\
GPT-oss-20B       & ${\sim}20$\,h & 78.3 & 96\,min & 3.2 \\
GPT-oss-120B      & ${\sim}23$\,h & 91.9 & 80\,min & 5.3 \\
\midrule
\textbf{Speedup}  & \multicolumn{2}{c}{} & \textbf{8--17$\times$} & \textbf{8--25$\times$} \\
\bottomrule
\end{tabular}
\end{table}

The cost gap has two compounding sources.
First, the $\mathcal{O}(N)$ call count: at a measured median of
$N{=}10$ step-level calls per TRAIL-GAIA trace and $N{=}30$ per
TRAIL-SWE-Bench trace, each \textsc{W2} trace consumes 10--30 model
invocations versus 1--2 for \textsc{+GI+SI} (the second pass of
\textsc{+GI+SI} fires only when Pass-1 detects a graph-source
category, on roughly 60--75\% of traces).
Second, the $\mathcal{O}(N^2)$ token volume: because each per-step
call carries the growing prefix of already-seen steps, the cumulative
input-token count is $\sum_{i=1}^{N}\mathcal{O}(i) =
\mathcal{O}(N^2)$, while \textsc{+GI+SI} prompts grow only as
$\mathcal{O}(N + |E|)$ where $|E|{=}13$ for the
intervention-validated causal subgraph.
The accuracy improvements reported for \textsc{+GI+SI} in
Table~\ref{tab:main_results} are therefore obtained at roughly
$\tfrac{1}{10}$th the wall time and $\tfrac{1}{20}$th the GPU-hours
of \textsc{W2}, on top of requiring none of \textsc{W2}'s
multi-window resubmission overhead.

\paragraph{Why we omit W3.}
The original \textsc{W3} is binary search: at each level the model
is asked to choose between the upper and lower half (exactly one),
recursing only into the chosen side until a leaf step is isolated.
This procedure produces \emph{one} prediction in $\mathcal{O}(\log
N)$ calls --- the entire reason \textsc{W3} is interesting in the
single-error setting.
TRAIL is multi-label: a label may occur in zero, one, or many spans,
and \emph{both} halves of any interval can independently contain
the target. A faithful multi-label adaptation must therefore (i)
run a separate bisection per category and (ii) allow both halves
to be positive at every level, recursing into all positive halves.
The first change multiplies the call count by $|\mathcal{C}|=19$;
the second eliminates the binary-search efficiency guarantee
because in the worst case (label present in every span) the
recursion enumerates all $N$ leaves with $\mathcal{O}(N)$
intermediate splits, giving $\mathcal{O}(|\mathcal{C}| \cdot N)$
calls --- strictly more expensive than \textsc{W2}'s
$\mathcal{O}(N)$ sweep, with no remaining algorithmic motivation.
We therefore include \textsc{W1} and \textsc{W2} as faithful
adaptations of the Who\&When framework, and exclude \textsc{W3}
on the grounds that its design hypothesis (sub-linear localization
of a unique error) is incompatible with TRAIL's multi-label
ground truth.

\paragraph{Implementation notes.}
All variants share the same JSON-extraction pipeline: a tolerant
parser that strips reasoning tags
(\texttt{<think>...</think>}/\texttt{<thinking>...</thinking>}),
removes the gpt-oss Harmony channel markers
(\texttt{assistantanalysis}/\texttt{commentary}/\texttt{final}),
and falls back to a brace-balanced object extractor when strict
\texttt{json.loads} fails. For reasoning models
(QwenLong-L1-32B, gpt-oss, DeepSeek-R1, QwQ) the runner
auto-bumps \texttt{max\_new\_tokens} from 8000 to 24000 because
the hidden \texttt{<think>} chain consumes most of the budget
before the answer block; without this lift the JSON is silently
truncated mid-string and every prediction collapses to empty.
A trailing rubric-scores call appends the five TRAIL evaluator
scores (\textit{reliability}, \textit{security},
\textit{instruction\_adherence}, \textit{plan\_optimization},
\textit{overall}) so the scorer can compute Pearson-$r$
correlations alongside W-F1, Loc, and Joint metrics.

\paragraph{Faithfulness audit: divergences from the original prompts.}
In prompting-based LLM-as-judge the \emph{prompt is the method}, so we
catalogue every textual change between the original prompts of
\citet{yin2025whoandwhen} (\texttt{Automated\_FA/Lib/utils.py}) and our
TRAIL adaptation. Each change is either \textbf{forced} (an unavoidable
consequence of the single$\to$multi remap and TRAIL's reference-free,
JSON-ingesting output schema) or a \textbf{deviation} (a choice that goes
beyond that remap). After the corrections described in this section,
\textsc{W1} and \textsc{W2} retain only forced changes plus a single
TRAIL-side adapter (the trailing rubric-scores call); \textsc{W3} is
excluded from headline experiments and listed here only to support that
exclusion.

\begin{itemize}\setlength{\itemsep}{0.4em}

\item \noindent\rule{\linewidth}{0.4pt}\\
      \textbf{Aspect:}\enspace \textsc{W1}/\textsc{W2}: TRAIL 19-category taxonomy block\\
      \textbf{Original W\&W:}\enspace absent (free-form ``identify the agent and the step'')\\
      \textbf{Our adaptation:}\enspace full leaf-category taxonomy injected at prompt head\\
      \textbf{Justified by single$\to$multi?:}\enspace \textbf{forced} --- categories are the multi-label output space\\
      \rule{\linewidth}{0.4pt}

\item \textbf{Aspect:}\enspace \textsc{W1}/\textsc{W2}: ground-truth answer (\texttt{ground\_truth})\\
      \textbf{Original W\&W:}\enspace included; the original authors annotate that it ``should be removed if ground truth shouldn't be used''\\
      \textbf{Our adaptation:}\enspace dropped\\
      \textbf{Justified by single$\to$multi?:}\enspace \textbf{forced} --- TRAIL evaluation is reference-free; including it would be data leakage\\
      \rule{\linewidth}{0.4pt}

\item \textbf{Aspect:}\enspace \textsc{W1}/\textsc{W2}: output format\\
      \textbf{Original W\&W:}\enspace free text (\texttt{Agent Name: \ldots Step Number: \ldots Reason: \ldots} or \texttt{1. Yes/No 2. Reason: \ldots})\\
      \textbf{Our adaptation:}\enspace strict JSON schema with an \texttt{errors} array of \texttt{(category, location, evidence, description, impact)} entries\\
      \textbf{Justified by single$\to$multi?:}\enspace \textbf{forced} --- multi-label output requires an array; the TRAIL scorer ingests JSON\\
      \rule{\linewidth}{0.4pt}

\item \textbf{Aspect:}\enspace \textsc{W2}: early exit on first ``yes''\\
      \textbf{Original W\&W:}\enspace defining behaviour --- the loop terminates as soon as one error is localized\\
      \textbf{Our adaptation:}\enspace removed; the sweep continues through every step-level span\\
      \textbf{Justified by single$\to$multi?:}\enspace \textbf{forced} --- a single-error early exit cannot recover the remaining \texttt{(category, span\_id)} pairs\\
      \rule{\linewidth}{0.4pt}

\item \textbf{Aspect:}\enspace \textsc{W1}/\textsc{W2}: trailing rubric-scores LLM call\\
      \textbf{Original W\&W:}\enspace absent (Who\&When predicts \texttt{(agent, step)} only; no rubric scores)\\
      \textbf{Our adaptation:}\enspace one extra zero-shot call per trace populates the five TRAIL evaluator scores\\
      \textbf{Justified by single$\to$multi?:}\enspace \textbf{deviation} (TRAIL-side adapter) --- needed for Pearson-$r$ correlations; isolated in a separate call so it does not contaminate the localization prompt\\
      \rule{\linewidth}{0.4pt}

\item \textbf{Aspect:}\enspace \textsc{W3}: scope of bisection\\
      \textbf{Original W\&W:}\enspace single combined run for the (single) error\\
      \textbf{Our adaptation:}\enspace 19 separate per-label runs (implemented but not used in headline experiments)\\
      \textbf{Justified by single$\to$multi?:}\enspace \textbf{forced} for multi-label, but multiplies cost/behaviour by $|\mathcal{C}|=19$\\
      \rule{\linewidth}{0.4pt}

\item \textbf{Aspect:}\enspace \textsc{W3}: half-selection at each level\\
      \textbf{Original W\&W:}\enspace pick exactly ONE of \{upper half, lower half\}\\
      \textbf{Our adaptation:}\enspace both halves can be marked positive; recurse into both\\
      \textbf{Justified by single$\to$multi?:}\enspace \textbf{deviation} --- algorithmically different from binary search; eliminates the $\mathcal{O}(\log N)$ guarantee and is the proximate reason \textsc{W3} is excluded from headline results\\
      \rule{\linewidth}{0.4pt}

\end{itemize}

The earlier draft of this baseline contained six additional deviations
(a span-index ``re-emphasize'' block in \textsc{W1}, a dropped task
description in \textsc{W1} and \textsc{W2}, a missing
``\emph{avoid being overly critical}'' calibration sentence in
\textsc{W2}, an assistant-prefix prefill of \texttt{\{"errors":} in
\textsc{W1}, a $K{=}20$ hard cap on the number of spans scanned by
\textsc{W2}, and a hard-coded \texttt{impact:"HIGH"} on every emitted
error rather than letting the model predict the
\{\textsc{HIGH},\textsc{MEDIUM},\textsc{LOW}\} severity as the canonical
TRAIL evaluator does). All six have been removed from the runner
(\texttt{baselines/who\_and\_when/run\_who\_and\_when\_vllm.py}); the
previous version is preserved at \texttt{old\_run\_who\_and\_when\_vllm.py}
for diffing. The remaining changes are forced by the multi-label
adaptation, with the trailing rubric-scores call as the only
non-forced TRAIL-side adapter.
